\typeout{NT FILE chap3.tex}%

\chapter{Requirements and Current System Analysis}
\label{cha:requirements_and_system_analysis}

\prependtographicspath{{5-Figures/Covers/}}

% epigraph configuration
\epigraphfontsize{\small\itshape}
\setlength\epigraphwidth{12.5cm}
\setlength\epigraphrule{0pt}

\section{Analysis of the Legacy System}
\label{sec:reqs_analysis_legacy}
\ntindex[Template]{}
    {
        The current internship management platform at \glsxtrshort{DI} at \glsxtrshort{NOVA} \glsxtrshort{FCT} is a web-based solution built on a deprecated technology stack. 
        Although it supports the undergraduate internship life cycle, it presents critical limitations regarding security, architecture, and functionality. 
        The section provides a systematic analysis of the current system's components and deficiencies, informed by direct exploration of the platform.
    }

    \subsection{Infrastructure and Tech Debt}
        {
            The system is hosted on a virtual machine and operates on an obsolete version of PHP dating back to approximately 2010, 
            but since the platform was made even before that maybe there are some modules that are way older. 
            This legacy codebase entails high maintenance risks and incompatibilities with modern libraries. 
            Specifically, older PHP versions no longer receive security updates, leaving the platform vulnerable to known exploits that malicious actors routinely target \cite{fastcomet-legacy-php-vulnerable}.
        }

        {
            A more critical vulnerability lies in the platform's network architecture: 
            the system communicates exclusively over \glsxtrshort{HTTP} (without \glsxtrshort{SSL}/\glsxtrshort{TLS} encryption), 
            representing a severe risk for transmitting sensitive student data and access credentials. 
            Data transmitted over unencrypted \glsxtrshort{HTTP} is exposed to interception and eavesdropping, allowing attackers to read login credentials, personal information, and institutional data in plaintext \cite{geeks4geeks-difference-http-https}.
        }

        {
            Email Deliverability Crisis: A particularly severe technical deficiency resides in the email infrastructure. 
            Due to the lack of modern authentication and deliverability configurations, transactional emails sent by the platform - intended for registration notifications, deadline reminders, call for proposals - are systematically flagged as spam by email providers \cite{smtp-improve-deliverability-transactional-mails}. 
            This breakdown in communication disrupts the workflow between the department, students, and companies, reducing the operational effectiveness of the platform regardless of other functionality.
        }

    \subsection{Functionalities and Limitations per Role}
        {
            Administration and Coordination: The administrator profile provides a global overview of entities and proposals, allowing actions such as bulk student imports via CSV and basic permission management. 
            A “\textit{superuser}” feature exists, allowing administrators to access any user account without credentials for troubleshooting purposes. However, a critical gap exists in audit logging: 
            while the system logs normal and superuser's logins and logouts, it lacks comprehensive action logging that tracks what users do during their sessions. 
            For compliance and accountability, the platform should record specific administrative actions such as proposal status changes, student data modifications, permissions adjustments, and bulk import operations.
        }

        {
            Content management relies on outdated, inflexible in-browser \glsxtrshort{WYSIWYG} editors for maintaining FAQs, Regulations, and Edition-specific information. 
            Mass communication is rigid: the system offers no configurable templates for different stages of the workflow, relying instead on a single standard message for inviting companies.
        }

        {
            Proposal Management and the Master thesis Gap: 
            The platform includes interfaces for companies or professors to submit internship proposals and Master’s dissertation themes. 
            However, these dissertation theme submission and management workflows are completely inoperative - they exist as “dead data” in the system, 
            forcing the department to manage dissertation assignments, publication and monitoring via external tools, negating the value of having collected this data in the platform.
        }

        {
            For undergraduate internships, the “Project Classification” workflow allows coordinators to accept, reject, or request clarifications by annotating proposals. 
            However, interaction remains limited: coordinators can toggle whether a proposal awaits company response, but the feedback mechanism lacks fluid integration for iterative proposal refinement.
        }

        {
            Students and Companies: The student interface allows ranking of internship preferences and the consultation of imported academic records. 
            For companies, human resources management is inefficient: there is no clear hierarchy or efficient mechanisms for managing multiple employees within a single company entity. 
            The overall \gls{UI} is outdated, hindering navigation and making it difficult for users to perceive current application status.
        }

    \subsection{Data and Reports}
        {
            The current system suffers from a critical lack of interoperability. 
            There are no native features to export complex data (student lists, proposals, results of attributions) into open, standard format such as \glsxtrshort{CSV} or \glsxtrshort{JSON}, which complicates operations that the coordinators might need to make with the data outside the platform. 
            Furthermore, company information is often outdated, as there are no periodic revalidation mechanisms to refresh company profiles and contact details ("Stale Company Data").
        }

\section{Requirements Gathering}
\label{sec:reqs_gathering}
\ntindex[Template]{}
    {
        Following the systematic analysis of the legacy platform and the identification of its critical failures, a comprehensive set of requirements was elicited for the new integrated platform.
        The elicitation process focused on addressing the needs of the diverse stakeholders involved in the academic lifecycle (students, companies, coordinators, secretariat, and administrators),
        with the goal of transitioning from fragmented spreadsheet-based management to a unified digital ecosystem supporting both undergraduate internships and master's dissertations.
        The resulting requirements were prioritized and refined in collaboration with the thesis advisers to ensure alignment with institutional constraints regarding data protection, interoperability, 
        and operation feasibility.
    }

    \subsection{Elicitation Methods}
        {
            To ensure the new platform addresses real operational pain points and does not introduce new barriers to existing workflows, 
            requirements were gathered through a multi-faceted approach:
        }
        \begin{itemize}
                \item {
                    \textbf{Legacy System Analysis}: Direct exploration and reverse-engineering of the existing platform to identify functional gaps (e.g. non-functional Master's modules) and technical debt (e.g. email deliverability issues).
                    This analysis, presented in Section \ref{sec:reqs_analysis_legacy}, served as the foundation for identifying core functional requirements.
                }
                \item {
                    \textbf{Document Analysis}: Review of current regulations, current internship rules, and manual processes currently used to manage Master's dissertations and feedback tracking. 
                    This provided insight into established business logic that must be preserved and systematized in the new platform.
                }
                \item {
                    \textbf{Workflow Modeling}: Mapping of the "\textit{as-is}" manual processes and system workflows against desired "\textit{to-be}" automated and digitized workflow.
                    Specifically for proposal submission and validation, student seriation algorithms, registrations workflows, company hierarchies, and feedback management.
                    Process models were documented to clarify state transitions, decision points, and inter-stakeholder communication patterns.
                }
                \item {
                    \textbf{Stakeholder Consultations}: Informal discussions with thesis advisers and the current coordinator responsible for the platform administration provided critical insights into administrative pain points,
                    compliance requirements, and feature priorities.
                    These discussions informed the prioritization of features and validated the proposed solution direction.
                }
        \end{itemize}

    \subsection{Stakeholders}
        {
            The system must support multiple distinct user groups. 
            However, the role hierarchy reflects the institutional structure where coordinators are designed professors with expanded administrative responsibilities,
            rather than a separate administrative role. 
        }
        \begin{itemize}
            \item{
                \textbf{Students}: Seek to consult and review available internship and dissertation/project proposals within their enrolled edition, 
                rank proposals according to their preferences, track the status of their candidacies and assignments, access their assigned supervisor's contact information, and, 
                if required by the coordinators, submit documentation (CV, motivation letter, or recommendations) for specific proposals.
            }
            \item{
                \textbf{Company Representatives}: Aim to register company information and manage organizational hierarchy (parent companies and sub-entities or departments),
                submit internship and Master's dissertations proposals during open call periods, receive and respond to feedback from coordinators regarding proposal clarifications, 
                manage company employees and designate contacts responsible for specific proposals, schedule and present their proposals during coordinator provided time slots, 
                and track interview progress and student candidacies through a dashboard. 
            }
            \item{
                \textbf{Professors}: Intent to propose scientific dissertation themes for the Master's degree, supervise students assigned to their dissertation topics,
                provide feedback on dissertation progress.
            }
            \item{
                \textbf{Edition Coordinators (Professor variant, Head Coordinator)}: Is a professor with comprehensive administrative responsibilities for the system. 
                Their responsibilities include configuring and managing editions (setting edition type - undergraduate internship or Master's dissertations, defining proposals' deadlines, 
                configuring available time slots for company presentations, creating edition-specific information pages such as FAQs and contact pages), 
                validate submitted proposals through a feedback-and-clarification workflow, control the attributions process (being able to manually add a student to an internship outside seriation),
                generating and customizing email templates for all communications stages, publish news and announcements targeting different user groups, managing user pre-registration and bulk imports via  \glsxtrshort{CSV},
                and overseeing the finalization of internship/dissertation protocols. 
                The head coordinator also maintains comprehensive audit logs and manages permissions assignments for sub-coordinators and other edition-specific administrative tasks.
            }
            \item{
                \textbf{Sub-Coordinators (Professor variant)}: Professors designed by the head coordinator to assist with proposals validation and feedback workflows for specific edition, 
                but without access to system-wide configuration. 
                They inherit permissions from the Professor's role but with additional privileges limited to proposal evaluation and validation, configuration of edition-specific resources (FAQs, contact pages, email templates) within their assigned edition.
            }
            \item{
                \textbf{Secretariat/Administrative Staff}: Manage financial and legal formalization of internships and dissertations, specifically verifying payments from companies for protocol signing and managing storage and archival of signed agreements.
                They require a view-only access to status dashboards showing active internships and their completion status.
            }
            \item{
                \textbf{Administrators}: Handle the infrastructure-level operations that fall outside the coordinator's scope, such as system backups, database maintenance, server updates, monitoring system health,
                managing  \glsxtrshort{SSL} certificates and security configurations, responding to critical technical incidents, and consulting detailed audit logs for technical troubleshooting and compliance verification.
                They do not manage user accounts, permissions, or edition-specific configurations - those remain the responsibility of the head coordinator. This role is typically filled by dedicated IT staff rather than a professor.
            }
        \end{itemize}

    \subsection{Functional Requirements}
        {
            The functional requirements are categorized by feature set to ensure modular development, 
            testing and clear traceability of stakeholder needs. Detailed specification for all functional requirements are presented in Annex \ref{ann:fr_nfr_annex}, Tables \ref{tab:freq-edition-config}-\ref{tab:freq-masters}. 
            The following subsections provide a summary of each category:
        }

        \begin{description}
            \item[FR-1x: Edition and Configuration Management:] 
                {
                    This category covers the administrative features coordinators need to create and manage academic editions.
                    It includes creating editions for undergraduate internships or Master's dissertations, configuring key timelines
                    (proposal submission, interviews, finalization), and managing edition-specific content such as FAQs, regulations and contact
                    information. 
                    Coordinators can also schedule company presentation slots and publish targeted announcements using reusable email templates,
                    replacing the legacy platform's rigid configuration mechanisms and enabling each edition to be tailored to its specific needs.

                }
            \item[FR-2x: Proposal Submission and Classification:] 
                {
                    This category groups the workflows through which companies and professors submit and refine internship or dissertation proposals.
                    The system supports detailed proposal submission, iterative clarification and feedback (rather than simple accept/reject), and final Classification
                    by coordinators. 
                    Companies can reuse proposals from previous editions, adapt them to new contexts, and schedule oral presentations in coordinator-defined time slots.
                    Automated notifications keep proposers informed about submission deadlines and feedback windows.
                }
            \item[FR-3x: Student Candidacies and Seriation:] 
                {
                    This category automates the process by which students apply to proposals and are ranked for assignment.
                    Students can browse available proposals for their edition, rank a limited number of preferences, and, when required,
                    attach supporting documents such as CVs or motivation letters.
                    The system executes a seriation algorithm based on predefined criteria (\textit{e.g.}, \glsxtrshort{ECTS}, average grade, coordinator-defined weights)
                    to produce ranked candidate lists for each proposal, notifies students of key milestones (ranking window, seriation results, confirmations), 
                    and allows companies to review and refine candidate selections, including excluding unsuitable candidates.
                }
            \item[FR-4x: Company and User Management:] 
                {
                    This category encompasses user and organizational management across all roles. 
                    Coordinators validate new company registrations and associated legal protocols before proposals can be submitted.
                    Company representatives manage their organization structure and employee accounts, keeping contacts and mentors for proposals up to date.
                    The system also supports bulk user pre-registering (\textit{e.g.}, student imports via CSV, manual registrations of professors and secretariat staff)
                    and a controlled "\textit{superuser}" capability, which allows coordinators to access user accounts for support purposes with full audit logging. 
                    Sub-coordinator roles can be assigned and revoked for specific editions to delegate proposal validation and feedback tasks.
                }
            \item[FR-5x: Communication and Notifications:] 
                {
                    This category covers the communication mechanisms that replace the legacy system's unreliable email workflows.
                    The platform automates the "\textit{Call for Proposals}", sending active companies edition-specific information 
                    (deadline, edition type) and allows coordinators to publish news and announcements targeted at specific user roles.
                    Critical events - such as new proposals, feedback requests, ranking deadlines and seriation results - generate both email
                    and in-app notifications, using configurable templates and a dedicated transactional email service to improve deliverability.
                    Coordinators can also send targeted email campaigns to selected user groups, with customizable subject and content.
                }
            \item[FR-6x: Reporting and Data Export:] 
                {
                    This category provides the reporting and audit capabilities required for institutional oversight.
                    Coordinators and administrator can export student lists, proposals, candidacies and final assignments in open formats (CSV, JSON),
                    overcoming the legacy platform's interoperability limitations.
                    The system maintains detailed audit log of critical actions (logins, proposal status, changes, permission updates, superuser access, bulk imports)
                    with timestamps and user attribution, as well as a trace of proposal feedback exchanges between coordinators and companies. 
                    These records support accountability, compliance checks and dispute resolution.
                }
            \item[FR-7x: Master's Dissertation Lifecycle Support:] 
                {
                    This category addresses the specific needs of Master's dissertations and engineering projects, turning previously unused "dead data" into fully supported workflows.
                    The system distinguishes between scientific dissertation and industry-oriented engineering projects, allowing professors and companies to propose themes, manage assigned students and
                    track progress.
                    Master's editions follow a lifecycle similar to undergraduate internships but include additional milestones such as topic approval and final submission/defense dates, reflecting the supervision
                    and evaluation requirements of Master's programmes
                }
        \end{description}

    \subsection{Non-Functional Requirements}
        {
            To address the technical debt of the legacy system and ensure the platform meets institutional and user expectations, 
            the new system must adhere to the following quality attributes, organized according to ISO/IEC 25010 standards. 
            Detailed specifications for all-non functional requirements are presented in Annex \ref{ann:fr_nfr_annex}, Tables \ref{tab:nfreq-security}-\ref{tab:nfreq-email-deliver}. 
            The following subsections provide a summary of each category:
        }
        \begin{description}
            \item[Security:] 
                {
                    Security requirements ensure that sensitive institutional and personal data is protected against unauthorized access.
                    The platform enforces \gls{RLS} at the database layer, so users only see data appropriate to their role and organization context,
                    and all communication is encrypted over \glsxtrshort{HTTPS}/\glsxtrshort{TLS} to avoid the legacy system's unencrypted \glsxtrshort{HTTP} risk \cite{geeks4geeks-difference-http-https}.
                    Passwords and credentials are stored using industry-standard hashing, and \gls{RBAC} with granular permissions limitations each
                    user to actions consistent with their assigned role \cite{pathlock-rbac-guide}. 
                }
            \item[Auditability \& Compliance:] 
                {
                   Auditability requirements guarantee that all critical actions can be traced for accountability, compliance and dispute resolution.
                   The system records a tamper-evident audit log of key events - such as logins, data changes, status updates, permission assignments and superuser sessions -
                   with precise timestamps and user attributions \cite{fortra-audit-log-best-practices}. 
                   Proposal feedback exchanges between coordinators are also logged, and superuser activity is recorded in detail (who accessed whom, when and for how long),
                   providing transparency over administrative interventions.
                }
            \item[Reliability \& Availability:] 
                {
                    This category ensures the platform remains operational during critical periods like proposal submission or seriation.
                    A transactional email service is used to deliver time-sensitive notifications reliably, addressing the legacy system's chronic email failures.
                    Automated backups of database, documents and configuration, combined with off-site replication, support recovery from data loss.
                    Health monitoring and alerting for components such as the database, backup jobs and email service enable fast detection and response to incidents, reducing downtime.    
                }
            \item[Performance \& Scalability:] 
                {
                    Performance requirements guarantee a responsive experience for users, even under peak load. 
                    User-facing pages such as dashboards, proposal lists and candidacy views are expected to load within a few seconds under normal conditions,
                    while key dashboards support real-time updates via WebSockets or server-sent events, so changes appear without manual refresh.
                    \glsxtrshort{API} endpoints must respond within defined latency targets for most requests, and rate limiting is applied to protect the system from abuse
                    and ensure fair use of resources across standard and bulk operations.
                }
            \item[Usability \& Accessibility:] 
                {
                    This category ensures the platform is easy to use and accessible, addressing the legacy system's poor and outdated user experience. 
                    The interface must be responsive on both desktop and mobile, with clear visual feedback on interactive elements (focus, hover, loading) so users understand the current state.
                    Core workflows - such as registration, proposal submission and candidacy preference ranking - should be completable without prior training, supported by concise inline
                    guidance and meaningful error messages.
                }
            \item[Interoperability \& Data Exchange:] 
                {
                    Interoperability requirements guarantee that data can flow reliably between the platform and other institutional tools.
                    The system supports bulk user imports (\textit{e.g.}, students via CSV, professors and secretariat staff) with validation, error reporting and rollback.
                    Coordinators and administrators can export key datasets - students, proposals, seriation results, assignments - in standard formats like
                    \glsxtrshort{CSV} or \glsxtrshort{JSON}, and a documented REST \glsxtrshort{API} allows third-party systems to query selected read-only data
                    without compromising security.
                }
            \item[Maintainability \& Extensibility:] 
                {
                    This category ensures the platform can be evolved and maintained without excessive technical debt.
                    The architecture is modular and layered (presentation, business logic, data access), allowing frontend, backend and database
                    concerns to be updated independently. 
                    The codebase is supported by appropriate documentation (README, API docs, comments for complex logic) and by disciplined Git workflows
                    with semantic versioning and meaningful commit history.
                    Configuration is externalized via environment variables or configuration files, rather than hardcoded, enabling consistent deployment across development,
                    staging and production environments.
                }
            \item[Email Deliverability (Critical Legacy Fix):] 
                {
                    Given that unreliable email was one of the most severe issues in the legacy system, email deliverability is treated as a separate non-functional requirement.
                    All outbound messages must use proper sender authentication to reduce spam classification and improve inbox placement.
                    The platform integrates with a professional transactional email service that offers monitoring and bounce handling, providing real-time insight into delivery status.
                    Email templates are tested before deployment to ensure correct rendering across common clients (Gmail, Outlook, Apple Mail) and on mobile devices.
                } 
        \end{description}

\section{Use Cases}
\label{sec:reqs_use_cases}
\ntindex[Template]{}
    {
        Use case modeling is a requirements engineering technique that describes interactions between actors (stakeholders) and the system to achieve specific goals.
        This section documents key use cases for the dissertation/internship management platform, representing critical workflows identified in the requirements' analysis. 
        Each use case follows a standardized format including actors, preconditions, postconditions, main flow, and alternative flows.
    }

    \subsection{Actors and Their Goals}
        {
            Bases on the stakeholder analysis in Section 3.2.2, the platform involves the following actors:
        }

        {
            \textbf{Primary Actors} (initiate use cases to achieve their own goals):
        }
        {
            \begin{itemize}
                \item {\textit{Student} - Seeks to discover proposals, rank preferences, and track candidacy status}
                \item {\textit{Company Representative} - Aims to submit proposals and manage company participation}
                \item {\textit{Professor} - Intends to propose dissertation themes and supervise students}
                \item {\textit{Edition Coordinator} - Responsible for configuring  editions and overseeing workflows}
            \end{itemize}
        }

        {
            Secondary Actors (support the system but do not initiate the primary goal):
        }
        {
            \begin{itemize}
                \item {\textit{Secretariat} - Manage financial and legal documentation}
                \item {\textit{System Administrator} - Maintains infrastructure and system health}
            \end{itemize}
        }

    \newpage
    \subsection{High-Level Use Case Diagram}
        \begin{figure}[htbp]
            \centering
            \includegraphics[height=7in]{MyFigures/use_case_high_level_diagram}
            \caption{High-Level Use Case Diagram of the system with core functionalities}
            \label{fig:Figure_Use_Case_Diagram_HighLevel}
        \end{figure}

    \subsection{Detailed Use Case Specifications}
        {
            This subsection documents the detailed specification of key use cases using the standard Cockburn template format:
            Use Case Name, Primary Actor(s), Secondary Actor(s), Preconditions, Postconditions, Main Flow, Alternative Flows \cite{visual-pardigm-uml-use-case-modeling}.
            The complete specifications for all use cases identified in this section are presented in Annex 2, Tables II.1-II.12.
        }

        {
            The following describes the scope and organization of documented use cases:
        }

        {
            Core Workflow Use Cases:
        }
        \begin{itemize}
            \item {
                \textbf{UC-01: Student Ranks Internship/Dissertation Preferences} - Central to the student candidacy phase;
                demonstrates the core functionality of the platform for students
            }
            \item {
                \textbf{UC-02: Coordinator Validates and Provides Feedback on Proposal} - Critical workflow; illustrates the 
                iterative refinement loop between coordinator and company
            }
            \item {
                \textbf{UC-03: System Executes Student Seriation Algorithm} - Core business logic; represents the automated 
                ranking that differentiates this platform form manual processes
            }
            \item {
                \textbf{UC-04: Company Responds to Coordinator Feedback and Resubmits Proposal} - Demonstrates bidirectional
                communication; shows how feedback loops improve proposal quality 
            }
            \item {
                \textbf{UC-05: Edition Coordinator Creates New Edition} - Administrative setup; establishes the foundation for
                almost all the other workflows
            }
            \item {
                \textbf{UC-06: Company Representatives Registers and Submits Initial Proposal} - Entry point for external stakeholders;
                illustrates onboarding and validation workflows.
            }
            \item {
                \textbf{UC-07: Student Submits Support Documentation} - Optional workflow, triggered when a coordinator requires CV or recommendation letters 
            }
        \end{itemize}
        
        {
            Secondary Workflow Use Cases:
        }
        \begin{itemize}
            \item {
                \textbf{UC-08: Company Views Assigned Candidates List and Schedule Interviews} - 
                Post seriation activities; enables company participation in selection
            }
            \item {
                \textbf{UC-09: Professor Proposes Dissertation Theme} - Master's degree specific; 
                operationalizes previously non-functional dissertation workflow 
            }
            \item {
                \textbf{UC-10: Secretariat Verifies Payment and Finalizes Protocol} - Post-assignment; ensures legal and 
                financial closure
            }
            \item {
                \textbf{UC-11: Edition Coordinator Publishes Call for Proposals} - Communication workflow; 
                ensures all companies receive the invite to send proposals to the edition
            }
            \item {
                \textbf{UC-12: Administrator Performs Superuser Access for Student Support} - Administrative support;
                includes detailed audit logging requirements
            }
        \end{itemize}

        {
            Each use case specification includes:
        }
        \begin{itemize}
            \item {Actors - Primary (initiates the goal) and secondary (supports the goal) participants}
            \item {Preconditions - System state and permissions required before the use case can begin}
            \item {Postconditions - Guaranteed system state after the use case completes (success of failure)}
            \item {Main Flow - The primary, "happy path" sequence of steps}
            \item {Alternative Flows - Exception handling, edge cases, and variant paths 
                (numbered A1, A2, A3, etc.) describing specific conditions where the main flow diverges.
            }
        \end{itemize}
        {
            Traceability to Requirements: Steps in the use case flows explicitly reference, when they have one, the corresponding Functional Requirement or Non-Functional Requirement (NFR-XX), 
            ensuring complete traceability between user interactions and system capabilities.
        }

    \subsection{Use Case Relationships and Extensions}
        {
            The use case model also captures a small set of relationships that express shared or conditional behavior between use cases.
        }

        {
            First, there is one extend relationship:
        }
            \begin{itemize}
                \item {
                    Respond to Feedback extends Submit Proposals when the coordinator has requested clarifications. In this situation, the company follows an additional flow
                    where it edits and resubmits the proposal instead of creating a completely new one.
                }
            \end{itemize}
        
        {
            Then, there are several include relationships, where one use case always performs another pas part of its normal execution:
        }
            \begin{itemize}
                \item {
                    Validate Proposals includes Send Feedback to Company, because every validation ends with the coordinator communicating 
                    an acceptance, rejection, or clarification request.
                }
                \item {
                    Execute Seriation includes System Generate Audit Log, since each seriation run must be recorded with its parameters and results
                    for auditability.
                }
                \item {
                    Publish Call for Proposals, Execute Seriation, and Validate Proposals each include System Notifies Company, reflecting that companies
                    are automatically informed when a new edition is open for proposals, when seriation results are available, or when their proposals change status.
                }

            \end{itemize}

    \subsection{Summary of Use Cases}
        {
            The use case model encompasses 12 primary use cases organized by the four main actors: 
            Students, Companies, Coordinators, and supporting roles (Professors, Secretariat)
        }

        {
            Student-Centric activities
        }
            {
                Students interact with the platform to browse and rank proposals, submit supporting documents if required,
                and track their candidacy status:
            }
            \begin{itemize}
                \item {UC-01: Rank Internship Preferences [FR-31]}
                \item {UC-07: Submit Supporting Documentation [FR-32]}
                \item {Track candidacy status after seriation [FR-34]}
            \end{itemize}
        
        {
            Company-Centric Activities
        }
            {
                Companies register, submit proposals, respond to coordinator feedback, and view ranked candidates:
            }
            \begin{itemize}
                \item {UC-06: Register and Submit Initial Proposal [FR-21, FR-41]}
                \item {UC-04: Respond to Coordinator Feedback and Resubmit Proposal [FR-22, FR-24]}
                \item {UC-08: View Ranked Candidates and Schedule Interviews [FR-35, FR-36]}
            \end{itemize}
        
        {
            Coordinator-Centric Activities
        }
            {
                Coordinators manage editions, validate proposals, publish calls for proposals, and execute seriation:
            }
            \begin{itemize}
                \item {UC-05: Create New Edition [FR-11, FR-12]}
                \item {UC-02: Validate Proposals [FR-22, FR-23]}
                \item {UC-03: Execute Seriation Algorithm [FR-33]}
            \end{itemize}
    
        {
            Professor-Centric Activities
        }
            {
                Professors primarily propose dissertation themes and supervise students.
            }
            \begin{itemize}
                \item {UC-09: Propose Dissertation Theme [FR-71, FR-72]}
            \end{itemize}

        {
            Cross-Functional Activities
        }
            {
                Supporting workflows for academic supervision, financial closure, and administrative support:
            }
            \begin{itemize}
                \item {UC-10: Verify Payments and Finalize Protocol [FR-41]}
                \item {UC-12: Superuser Access for Student Support [FR-44, NFR-AUD-4]}
            \end{itemize}

        {
            Conclusion
        }
        
        {
            Complete detailed specifications for all 12 use cases - including preconditions, postconditions, main flows, 
            and alternative flows - are documented in Annex 2, Tables II.1-II.12. 
            These specifications will serve as input for the system architecture design and form the basis for implementation and testing.
        }
